%% IEEE conference manuscript — v2-campaign methodology.
%% Faithful LaTeX transcription of ../artigo_ieee_v2_en.md.
%% PENDING markers: numbers/claims to fill from the v2 campaign
%% (results_controlled/heuristics_metaheuristics_v2/), rendered in red via
%% \pending{...}. Toggle \draftpendingtrue/false to show/hide them.
\documentclass[conference]{IEEEtran}
\IEEEoverridecommandlockouts

\usepackage{cite}
\usepackage{amsmath,amssymb,amsfonts}
\usepackage{algorithmic}
\usepackage{graphicx}
\usepackage{textcomp}
\usepackage{xcolor}
\usepackage{booktabs}
\usepackage{array}

% --- PENDING (draft) marker ------------------------------------------------
\newif\ifdraftpending
\draftpendingtrue  % set \draftpendingfalse for a clean (marker-free) build
\newcommand{\pending}[1]{%
  \ifdraftpending{\color{red}\textbf{[PENDING-v2: #1]}}\else{}\fi}

\begin{document}

\title{Constraint-Aware Metaheuristic Optimization of\\
Slice Allocation Weights in 5G/O-RAN Networks}

\author{\IEEEauthorblockN{\pending{author names}}
\IEEEauthorblockA{\pending{affiliation, city, country}\\
\pending{emails}}}

\maketitle

\begin{abstract}
Network slicing requires radio access networks to partition resources among
service classes with conflicting requirements: throughput for eMBB, latency
and reliability for URLLC, and massive connectivity for MTC. We formulate
inter-slice resource allocation as a continuous optimization problem over a
weight vector on the probability simplex, in which each weight controls one
slice's share of the physical resource block (PRB) budget, while conventional
schedulers (proportional fair) remain in charge of intra-slice allocation. We
compare four metaheuristics---Particle Swarm Optimization (PSO), a Genetic
Algorithm (GA), Simulated Annealing (SA), and a hybrid method combining
population-based exploration with stochastic local refinement---under a
\emph{feasibility-first objective} that treats slice service-level agreements
(SLAs) as hard constraints: URLLC packet-level 99th-percentile delay,
per-slice packet delivery ratio, and eMBB offered-load satisfaction. Unlike
weighted-sum formulations, this objective cannot reward configurations that
starve a slice. Each candidate is evaluated by a full-stack 5G simulation
(ns-3 / 5G-LENA) with per-packet latency instrumentation, across
\pending{4} traffic scenarios and 10 independent seeds, under an equalized
evaluation budget of 48 evaluations per method. Paired non-parametric tests
(Friedman, Wilcoxon signed-rank) over per-seed best scores show that
\pending{main statistical finding}. Under congestion, the best evolved weight
vectors \pending{headline physical result: deadline-violation rate / worst-slice
SLA vs.\ baselines}. All code, scenarios, and per-seed artifacts are available
for reproduction.
\end{abstract}

\begin{IEEEkeywords}
network slicing, 5G, O-RAN, metaheuristics, constrained optimization, particle
swarm optimization, genetic algorithms, simulated annealing, resource
allocation.
\end{IEEEkeywords}

\section{Introduction}
Fifth-generation (5G) networks support network slicing natively, allowing
heterogeneous service classes to share one physical infrastructure under
distinct service-level agreements (SLAs). Enhanced Mobile Broadband (eMBB)
slices prioritize sustained throughput; Ultra-Reliable Low-Latency
Communication (URLLC) slices prioritize bounded latency and delivery
reliability; massive Machine-Type Communication (MTC) slices prioritize
serving many low-rate devices without starvation. Because these objectives
conflict whenever radio resources are scarce, inter-slice allocation is a
multi-objective decision problem whose difficulty grows with load.

O-RAN architectures encourage programmable control loops (rApps/xApps) that
adjust radio-layer policy without replacing the underlying schedulers. A
pragmatic design keeps a standard intra-slice scheduler (e.g., proportional
fair) and optimizes, at the inter-slice level, a small continuous vector of
allocation weights that partitions the PRB budget among slices. This
separation retains compatibility with production schedulers and reduces the
search space to a low-dimensional simplex---but the mapping from weights to
network KPIs remains non-convex, noisy, and only accessible through simulation
or measurement, which motivates black-box metaheuristic search.

A central methodological problem, often overlooked, is the \emph{objective
function}. Weighted-sum scores over KPI aggregates can assign high scores to
degenerate allocations: in a preliminary version of this study, a weighted-sum
objective ranked \emph{first} a configuration that effectively starved the MTC
slice (packet delivery ratio of 2.4\%, mean delay above 4\,s) because a
throughput-only SLA indicator saturated at 100\%. This is not an implementation
accident but a structural failure mode of soft-penalty scalarizations. We
therefore adopt a feasibility-first formulation in which SLA violations are
hard constraints and infeasible candidates are always dominated by feasible
ones, and we instrument the simulator with \emph{per-packet} latency
percentiles so that tail-latency constraints are measured rather than
approximated by window averages.

The contributions of this paper are:
\begin{enumerate}
\item A constraint-aware formulation of inter-slice weight optimization on the
simplex, with hard SLA constraints on URLLC tail latency (per-packet p99),
per-slice delivery ratio, and eMBB offered-load satisfaction.
\item A controlled comparison of PSO, GA, SA, and a hybrid metaheuristic under
an equalized evaluation budget (48 evaluations per method) across \pending{4}
traffic scenarios and 10 independent seeds, with paired non-parametric
statistical testing.
\item A per-packet measurement methodology for slice SLAs in ns-3/5G-LENA---
packet-level delay percentiles (p95/p99/p99.9), deadline-violation rate, and
in-time reliability---replacing window-sampled statistics that systematically
underestimate tails.
\item A physical-KPI evaluation of the best evolved allocations against pure
schedulers (RR, PF, BCQI), slice-aware baselines, and adaptive heuristics
(including AQPS), showing when weight optimization pays off and when it does
not.
\end{enumerate}

\section{Related Work}
\pending{related-work section: RSLAQ (DRL-based SLA control in O-RAN)~\cite{rslaq},
AQPS (QoS-aware priority scheduling)~\cite{aqps}, metaheuristics for RAN
resource allocation, network-slicing surveys. Complete after reference
consolidation.}

Distinctly from DRL approaches such as RSLAQ~\cite{rslaq}, which learn a control
policy online, our approach searches offline for a static per-scenario weight
configuration; it requires no training infrastructure, is auditable (every
evaluated candidate is persisted with its full KPI record), and its outputs can
seed either static configuration (Non-RT RIC time scale) or serve as ground
truth for learned policies. It is \emph{not} a Near-RT control mechanism: each
candidate evaluation costs a full simulation, which confines metaheuristic
search to offline/periodic optimization.

\section{Problem Formulation}
\subsection{Decision variables}
Let $S = 3$ slices indexed by $i \in \{\mathrm{eMBB}, \mathrm{URLLC},
\mathrm{MTC}\}$. A candidate solution is the continuous weight vector
\begin{equation}
\mathbf{w} = [w_{\mathrm{eMBB}}, w_{\mathrm{URLLC}}, w_{\mathrm{MTC}}], \quad
w_i \ge 0, \quad \sum_{i} w_i = 1,
\end{equation}
i.e., a point on the 2-simplex. The MAC scheduler partitions the per-slot
resource block group (RBG) budget proportionally to $\mathbf{w}$ among slices
with active demand (budget unused by inactive slices is redistributed), and a
proportional-fair policy allocates within each slice. Candidates produced by
the optimizers are projected back onto the simplex by clipping negative
components and renormalizing.

\subsection{Constraint-aware objective}
Each candidate is evaluated by one end-to-end simulation, from which per-slice
KPIs are extracted. The objective is \emph{feasibility-first} maximization:
\begin{equation}
F(\mathbf{w}) =
\begin{cases}
\dfrac{100}{S}\sum_{i} \mathrm{SLA}_i(\mathbf{w}) & \text{if } V(\mathbf{w}) = 0,\\[2mm]
-100\, V(\mathbf{w}) & \text{otherwise,}
\end{cases}
\end{equation}
where $V(\mathbf{w}) \ge 0$ is the sum of normalized violations of the hard
constraints
\begin{align}
D^{99}_{\mathrm{URLLC}} &\le 10~\mathrm{ms}, &
\mathrm{PDR}_{\mathrm{URLLC}} &\ge P^{\min}_{U}, \\
\mathrm{PDR}_{\mathrm{MTC}} &\ge P^{\min}_{M}, &
T_{\mathrm{eMBB}} &\ge \phi\, L_{\mathrm{eMBB}},
\end{align}
with $D^{99}$ the \emph{per-packet} 99th-percentile downlink delay, PDR the
end-to-end packet delivery ratio, $T$ the delivered throughput, $L$ the offered
load, and $\phi$ the minimum served fraction. $\mathrm{SLA}_i \in [0,1]$ is a
composite per-slice satisfaction, $\mathrm{SLA}_i = \min(S^{thr}_i, S^{pdr}_i,
S^{delay}_i)$---the \emph{worst} dimension governs, so no slice dimension can
compensate another. Two properties follow by construction: (i) any feasible
candidate dominates every infeasible one; (ii) among infeasible candidates the
search still receives a gradient toward feasibility, avoiding flat penalty
regions.

Constraint thresholds were calibrated on a pilot campaign so that every
scenario retains a non-empty feasible region:
\pending{final calibrated $P^{\min}_U$, $P^{\min}_M$, $\phi$; pre-calibration
defaults 90\%, 80\%, 0.60. Report the calibration outcome here.}

This formulation operationalizes the URLLC latency budget as a measured
packet-level constraint rather than a soft penalty on mean delay---a design
decision motivated by the failure mode described in Section~\ref{sec:intro}.

\section{Optimization Methods}
All methods maximize the same $F(\mathbf{w})$, receive the same evaluation
budget, and share the checkpointed evaluation pipeline (Section~\ref{sec:campaign}).

\textbf{PSO.} A swarm of 6 particles with inertia 0.55 and cognitive/social
coefficients 1.30; positions are renormalized onto the simplex after each
update. $8 \times 6 = 48$ evaluations.

\textbf{GA.} Population 6 with elitism, blend crossover on the simplex, and
Gaussian mutation ($\sigma = 0.12$) followed by renormalization. $8 \times 6 =
48$ evaluations.

\textbf{SA.} Single-trajectory annealing with Gaussian perturbations and a
decreasing temperature schedule. To equalize budgets, the chain length is set
to 48 evaluations (matching GA/PSO), not to the iteration count.

\textbf{Hybrid.} A population method combining an elite archive, PSO-inspired
updates, GA-style recombination, and one SA-style local refinement of the
global best per iteration. Because of the extra local step it consumes 56
evaluations ($8 \times (6{+}1)$); comparisons at the common 48-evaluation
ceiling are also reported to keep the budget comparison fair.
\pending{verify both budget readings appear in the results tables.}

The optimizer RNG is seeded independently per simulation seed
($\mathit{search\_seed} = \mathit{base\_seed} + \mathit{ns3\_seed}$), so the 10
repetitions are independent searches, not one trajectory replayed on 10
channels.

\section{Experimental Methodology}\label{sec:campaign}
\subsection{Simulation environment}
Experiments use ns-3 with the 5G-LENA NR module~\cite{5glena}. A single gNB
(3.55~GHz, 100~MHz, numerology $\mu{=}1$, TDD \texttt{D|D|8D|4GB|4U|U|U}, RLC
UM, 43~dBm) serves stationary UEs placed on a disk of radius 100~m; each UE
receives a downlink UDP flow from a remote host through the core network. UEs
are grouped into the three slices in contiguous ID ranges; FlowMonitor maps
flows back to slices by destination port. A custom slice-aware MAC scheduler
partitions the RBG budget according to $\mathbf{w}$ (Section~III-A).

\subsection{Traffic scenarios}
\pending{confirm final scenario table. Four scenarios in the main paper
(low\_traffic, normal, congestion, stressed); a fifth profile
(insufficient\_resources) is simulated and reported as a robustness check.}

\begin{table}[t]
\centering
\caption{Traffic scenarios (nominal offered load).}
\label{tab:scenarios}
\begin{tabular}{lccc}
\toprule
Scenario & UEs (eMBB/URLLC/MTC) & Offered load (Mb/s) \\
\midrule
low\_traffic & 2 / 2 / 6 & 57 \\
normal & 5 / 5 / 10 & 73 \\
stressed & \pending{n} & 128 \\
congestion & 15 / 10 / 35 & 245 \\
\bottomrule
\end{tabular}
\end{table}

These are nominal input profiles; the operational regime (whether the cell is
actually congested) is diagnosed from measured PDR and delay, not assumed from
the label. Simulation time is 15~s per evaluation with 0.4~s application
warm-up and 0.2~s drain (KPIs are computed over the active window only).

\subsection{Per-packet SLA instrumentation}
Latency statistics are computed from the FlowMonitor \emph{per-packet delay
histogram} (bin width 0.5~ms), merged per slice: percentiles p95/p99/p99.9 are
interpolated within bins; the \emph{deadline-violation rate} is the fraction of
delivered packets exceeding the slice's latency budget; \emph{in-time
reliability} is the fraction of transmitted packets delivered within budget
(undelivered packets count as violations). Packet loss is reported end to end
($\mathrm{PLR} = (tx - rx)/tx$, so $\mathrm{PDR} + \mathrm{PLR} = 100\%$ by
construction), with detected-loss (FlowMonitor) kept as a separate diagnostic.
We emphasize this because window-averaged sampling---used in a preliminary
version of this pipeline---produced $\text{mean} > \text{p95}$ artifacts in
about one fifth of the measured cells and systematically hid tail latency; all
results in this paper use per-packet statistics.

\subsection{Campaign protocol and reproducibility}
For every (scenario, seed) pair, all four methods run under the same objective,
the same seeds (1--10), and the same simulator build. Every candidate
evaluation is persisted with a JSON sidecar (weights, score, objective version,
full KPI record), which makes the search resumable and the comparison auditable
a posteriori. Baselines (pure RR/PF/BCQI, slice-aware variants, adaptive
heuristics including AQPS) run under the identical protocol and seeds.
\pending{total evaluation counts and total compute time of the campaign.}

\subsection{Statistical analysis}
The experimental unit is the \emph{best score per (method, scenario, seed)}---
one independent search outcome per seed, 10 per method/scenario. We report
mean\,$\pm$\,sd, median, and bootstrap 95\% CIs; the omnibus Friedman test
across the four methods; and pairwise Wilcoxon signed-rank tests with
rank-biserial effect sizes~\cite{demsar}. With $n = 10$ paired seeds these
tests are adequately powered for large effects; we do not claim significance
where $p \ge 0.05$.

\section{Results}
All numbers below come exclusively from the v2 campaign. Tables are filled by
\texttt{analyze\_v2\_stats.py}; figures by \texttt{generate\_v2\_audit\_figures.py}.

\subsection{Feasibility structure of the search space}
\pending{per scenario: fraction of evaluated candidates satisfying all hard
constraints; which constraint binds most often (expected: URLLC p99 under
congestion); whether any scenario has an empty feasible region after
calibration.}

\subsection{Method comparison under equal budget}
\pending{Table: best score per seed (mean$\pm$sd, median, CI95) per
method$\times$scenario; Friedman $\chi^2$/p per scenario; pairwise Wilcoxon
$p$ + effect sizes. Figures: per-seed dispersion strips; score bars.}

\pending{Convergence: evaluations to reach 95\% of final best, per method.}

\subsection{Physical SLA outcomes of the best allocations}
\pending{Table/figures: URLLC p99 and p99.9, deadline-violation rate,
worst-slice SLA, MTC PDR per method. Key question: does ANY method satisfy all
SLAs simultaneously, and at what load does that stop being possible?}

\subsection{Comparison with schedulers and heuristics}
\pending{best metaheuristic vs pure RR/PF/BCQI, slice-aware baselines, AQPS and
adaptive heuristics---same protocol, same seeds. Physical Pareto: total
throughput $\times$ URLLC p99.}

\subsection{Weight stability across seeds}
\pending{violin of best-candidate weights per slice across seeds; coefficient
of variation per scenario---whether the optimum is a property of the scenario
or of the noise realization.}

\section{Discussion}
\pending{rewrite around actual findings. Anchor points:}
\begin{enumerate}
\item \emph{When does weight optimization pay off?} Preliminary evidence
indicates gains concentrate under resource contention; quantify with v2.
\item \emph{Feasibility vs.\ score.} Document how the recommended allocations
differ from the weighted-sum optima (which starved MTC).
\item \emph{Practical deployment.} Search cost per scenario
(\pending{CPU-hours}) confines the approach to Non-RT RIC time scales; discuss
the xApp path and its relation to learned policies.
\item \emph{Method choice.} \pending{whether any method is statistically
distinguishable at $n{=}10$; if not, say so plainly---a negative result with
adequate testing is publishable and honest.}
\end{enumerate}

\section{Threats to Validity}
\begin{enumerate}
\item \textbf{Simulation scope.} Single cell, stationary UEs, downlink UDP
only; topology and load co-vary across scenarios.
\item \textbf{Statistical power.} $n = 10$ paired seeds supports detection of
large effects only; small differences may remain undetected. We report effect
sizes and CIs alongside p-values.
\item \textbf{Budget asymmetry.} The hybrid consumes 56 evaluations vs.\ 48;
both full-budget and common-ceiling comparisons are reported.
\item \textbf{Threshold calibration.} Hard-constraint thresholds were
calibrated on a pilot and held fixed; different SLA contracts would change the
feasible region. The calibration procedure and sensitivity are documented.
\item \textbf{Evaluation noise.} A single simulation evaluates each candidate;
seed effects are handled across, not within, searches.
\pending{report any evaluation timeouts/crashes and their hard-penalty handling.}
\item \textbf{Measurement caveats.} FlowMonitor throughput includes IP/UDP
headers; delay histograms have 0.5~ms resolution (percentile error bounded by
bin width); histogram-based percentiles are interpolated.
\end{enumerate}

\section{Conclusion}
\pending{write after results. Required elements: (i) the constraint-aware
formulation and per-packet instrumentation as the methodological contribution;
(ii) the statistically-tested method comparison (state clearly if methods are
indistinguishable); (iii) the physical-KPI verdict---does optimized slicing
satisfy all SLAs simultaneously, and under which loads; (iv) future work:
online adaptation (xApp), multi-cell, mobility, learned policies distilled from
the search.}

\bibliographystyle{IEEEtran}
\bibliography{references}

\end{document}
